%% =====================================================================
%% part_viii_synthesis.tex
%%
%% Vol II Part VIII overview: "From Frontier to Theorem"
%%
%% Synthesis-only file. Cross-references existing theorems by \label;
%% contains NO new theorem statements and NO duplicated proofs. Each
%% section is a ~150-300 word narrative mapping one of the four closed
%% frontiers (or one supporting structural result) into the
%% programme's Platonic form.
%%
%% Location: chapters/frame/ (Vol II).
%% Inputs: chapters/theory/curved_dunn_higher_genus.tex
%% chapters/theory/chiral_higher_deligne.tex
%% chapters/theory/koszulness_moduli_M_kosz.tex
%% chapters/theory/infinite_fingerprint_classification.tex
%% chapters/connections/brace.tex
%% chapters/connections/bv_brst.tex
%% (Vol I) chapters/theory/periodic_cdg_admissible.tex
%% =====================================================================

%%%%%%%%%%%%%%%%%%%%%%%%%%%%%%%%%%%%%%%%%%%%%%%%%%%%%%%%%%%%%%%%%%%%%%%%%%
\chapter*{Part VIII Synthesis: Theorem Atlas}
\addcontentsline{toc}{chapter}{Part VIII Synthesis: Theorem Atlas}
\label{chap:part-viii-synthesis}
%%%%%%%%%%%%%%%%%%%%%%%%%%%%%%%%%%%%%%%%%%%%%%%%%%%%%%%%%%%%%%%%%%%%%%%%%%

%% =====================================
\section{Theorem atlas: overview}
\label{sec:part-viii-overview}

Throughout Parts I--VII, $\SCchtop$ remains two-coloured until the
inner-stress topologisation data is supplied; on those loci the
$E_{1}$-chiral seed acquires the topological rungs leading to the
$E_{3}^{\mathrm{top}}$ climax. The theorem locus of that ladder is cut out by
four hypotheses:
(a) modular-bootstrap \(H^{2}\)-acyclicity plus the degree-\(2\)
comparison for curved-Dunn transport at \(g\ge2\);
(b) DS--Hochschild chain-level compatibility for class~$\mathbf{M}$;
(c) periodic-CDG admissible-level Koszul compatibility for
Kazhdan--Lusztig; (d) chain-level chiral Deligne--Tamarkin after a
fixed Drinfeld-associator gauge.

The four threads converge on a single object: the Koszulness Moduli
Scheme $M_{\mathrm{Kosz}}(\mathcal A)$ with its $\mathrm{GRT}_{1}(\Q)$-torsor
of fourteen charts (Theorem~\ref{thm:mkosz-fourteen-charts}). Each of
(a)--(d) is a coordinate-chart identification in a fixed
Drinfeld-associator gauge. The four frontiers are not independent
obstructions; they are four shadows of one Platonic object.

%% =====================================
\section{Curved-Dunn $H^{2}=0$ as a theorem}
\label{sec:part-viii-curved-dunn}

Classical Dunn additivity reads
$E_{2} \simeq E_{1} \otimes_{\mathrm{Dunn}} E_{1}$. For the $\SCchtop$
genus tower the analogous statement demands a curved refinement: the
two colours tensor against a shared Arakelov background contributing
curvature \(d_{B}^{2}=\kappaChHodge(\mathcal A)\cdot \omega_g\) at each
genus. Curved Dunn additivity descends from genus~0 (trivial) and genus~1
(twisted tensor criterion) to $g \ge 2$ subject to the $H^{2}$ vanishing
of the curved-Dunn twisting-cochain complex. The genus-$1$ input is the
Gauss--Manin uncurving of
Proposition~\ref{prop:genus1-twisted-tensor-product}; the higher-genus
obstruction is the $H^{2}$ class.

For \(g\ge2\), the curved-Dunn twisting-cochain complex has \(H^{2}=0\)
under modular-bootstrap \(H^2\)-acyclicity and the degree-\(2\)
comparison (Theorem~\ref{thm:curved-dunn-H2-vanishing-all-genera}).
The bridge chain map of
Proposition~\ref{prop:modular-bootstrap-to-curved-dunn-bridge}
transports modular-bootstrap vanishing to the curved-Dunn complex.
Combined with the Jimbo--Miwa irregular-singular KZB regularity
(Theorem~\ref{thm:irregular-singular-kzb-regularity}) from a boundary
formal type, modular operad composition promotes from integral to
generic non-integral level at all $(g,n)$
(Corollary~\ref{cor:modular-operad-composition-all-levels}).

%% =====================================
\section{The modular bootstrap $\to$ curved-Dunn bridge chain map}
\label{sec:part-viii-bridge}

The load-bearing construction underlying
Section~\ref{sec:part-viii-curved-dunn} is the bridge chain map
$\Phi : \BMBC^{\bullet} \to \CDTC^{\bullet}$ of
Proposition~\ref{prop:modular-bootstrap-to-curved-dunn-bridge}.
The map promotes the per-genus vanishing $H^{2}(\BMBC) = 0$
(Definition~\ref{def:modular-bootstrap-complex}) to
$H^{2}(\CDTC) = 0$ in the curved-Dunn complex
(Definition~\ref{def:curved-dunn-complex}). Directionality is essential:
the curved-Dunn complex sees the tensor factorisation but not the
bootstrap constraints; the modular-bootstrap complex sees the
constraints but not the factorisation. $\Phi$ is the zigzag relating
them through a common gauge-fixed resolution.

Modular-bootstrap vanishing in degree \(2\) transports to the
curved-Dunn obstruction complex with no additional Dunn-additivity
input. The per-genus $H^{2}(\BMBC) = 0$ was visible classically through
Beilinson--Drinfeld chiral Koszul duality and the clutching structure
of the modular operad; the chain-level argument that the vanishing
survives the curved-Dunn tensor twist is the new content. The twisting
cochain for $\Phi$ is explicit in genus~1
(Proposition~\ref{prop:genus1-twisted-tensor-product}) and inductive in
$g$ via the stable-graph stratification of
$\overline{\mathcal{M}}_{g,n}$.

The bridge takes modular-bootstrap vanishing as input and emits
curved-Dunn vanishing as output; no circular appeal to Dunn additivity
enters the proof (Remark~\ref{rem:bridge-provenance}).

%% =====================================
\section{DS--Hochschild compatibility bridge for class $\mathbf{M}$}
\label{sec:part-viii-ds-hochschild}

For chiral algebras in classes G, L, C, the global triangle
$Z^{\mathrm{der}}_{\mathrm{ch}}(\mathcal{A}) \simeq
\mathcal{A}^{\mathrm{bulk}}$ identifying the derived chiral centre
with a 3d HT bulk was proved at chain level via boundary-linearity and
Costello--Gaiotto. For class~$\mathbf{M}$ (Virasoro, $W_{N}$, quartic-pole
lineage), chain-level identification demands a further theorem:
Drinfeld--Sokolov reduction transports $V_{k}(\mathfrak{g})$ to
$W_{k}(\mathfrak{g},f)$ at the level of chiral algebras, and chiral
Hochschild cohomology commutes with DS at the chain level.

The bridge is Theorem~\ref{thm:chd-ds-hochschild} of
Chapter~\ref{ch:chiral-higher-deligne}: the chain-level
quasi-isomorphism
$\mathrm{ChirHoch}^{\bullet}(W_{k}(\mathfrak{g},f)) \simeq
H^{\bullet}_{\mathrm{DS}}(\mathrm{ChirHoch}^{\bullet}(V_{k}(\mathfrak{g})))$.
The proof passes through three steps: (i) Homological Perturbation
Lemma transfer through the Drinfeld--Sokolov strong deformation
retract, strict weight-by-weight because the BRST filtration has
finite length on every weight space; (ii) Arakawa $C_{2}$-cofiniteness
lifting to chain-level finiteness; (iii) the HKR identification
$\mathrm{ChirHoch}^{\bullet}(V_{k}) \simeq
\mathcal{O}(T^{\ast}[-1]\,\DerM_{\mathrm{vac}}(V_{k}))$ followed by
BV reduction along the DS moment map.

Corollary~\ref{cor:universal-holography-class-M} packages this into
the universal holography statement for class~$\mathbf{M}$ in the
weight-completed/pro ambient: every conformal-vectored non-critical
class~$\mathbf{M}$ chiral algebra is the boundary of a canonical 3d HT
theory, with bulk its completed derived chiral centre. The strict
original-complex $E_3$ chain lift remains conditional as recorded in
the chiral Higher Deligne and universal-holography chapters.

%% =====================================
\section{Class $\mathbf{M}$ chain level from weight-completed HPL transfer}
\label{sec:part-viii-class-m-weight-completed}

The engine powering Theorem~\ref{thm:chd-ds-hochschild} is the
\emph{weight-completed} identification of BV, BRST, and bar structures
on class~$\mathbf{M}$. The raw direct-sum chain-level identification
fails: the direct-sum symplectic polarisation degenerates for
quartic-pole families because the pairing does not factor weight-by-weight.
The correct ambient is therefore different.

Proposition~\ref{prop:bv-bar-class-m-weight-completed} of
Chapter~\ref{ch:bv-brst} provides it: the weight-completed envelope
in which the BV--BRST--bar chain-level identification holds
unconditionally across all four shadow classes G, L, C, $\mathbf{M}$.
The completion is along the Li filtration by conformal weight; the
identification is strict on each weight space and extended to the
completion by the standard pro-nilpotent completeness argument.

The class~$\mathbf{M}$ side of the weight-completed identification
feeds Step~(ii) of the proof of Theorem~\ref{thm:chd-ds-hochschild}
via Arakawa $C_{2}$-cofiniteness. Weight-completeness permits HPL
transfer to commute with cohomological $\mathrm{ChirHoch}^{\bullet}$
on each conformal-weight stratum; Arakawa supplies the finiteness
control making the HPL argument converge in the limit. The chain-level
quasi-isomorphism survives completion, which is what permits
Corollary~\ref{cor:universal-holography-class-M} to be stated
unconditionally in the weight-completed/pro ambient on the
conformal-vectored non-critical locus.

%% =====================================
\section{Periodic-CDG at admissible Kac--Moody level}
\label{sec:part-viii-periodic-cdg}

The Kazhdan--Lusztig bar--cobar adjunction at admissible Kac--Moody
level demands periodicity of the category $\mathcal{O}$ under Koszul
dualisation on the semisimplified level — the periodic-CDG
Koszul-compatibility problem. Without it, the KL adjunction at
admissible $k$ has only generic-level force in downstream Vol~II
applications.

Theorem~\ref*{V1-thm:periodic-cdg-is-koszul-compatible}
(Vol~I Chapter~\ref*{V1-chap:periodic-cdg-admissible}) settles the
question: the periodic-CDG structure is Koszul-compatible on the
admissible locus. The proof invokes (i) Arakawa $C_{2}$-cofiniteness
at admissible level; (ii) Finkelberg tilting semisimplification;
(iii) Lusztig Tate periodicity of admissible blocks; (iv) commutation
of screening-adjoint with periodic-CDG, established by an independent
screening computation.

Theorem~\ref*{V1-thm:admissible-kl-bar-cobar-adjunction} then yields
the Kazhdan--Lusztig bar--cobar adjunction on the admissible KL locus
satisfying the exactness and periodic-CDG hypotheses. Every Vol~II
downstream theorem citing ``KL at generic $k$; admissible level
conditional on periodic-CDG'' may cite the admissible theorem directly
on that locus. The degenerate admissible logarithmic lane
(Creutzig--Kanade--Linshaw machinery) is treated separately in Vol~I.

%% =====================================
\section{Adams-analog construction for semisimplified Kazhdan--Lusztig}
\label{sec:part-viii-adams-analog}

The Adams spectral sequence analogy reads the bar complex as a
cohomological resolution whose ``Steenrod-algebra-analog'' primary
operations appear as $d_{1}$ differentials on the $E_{1}$-page. Two
questions remain: which chiral endomorphism functor plays the role of
the mod-$p$ Steenrod algebra; whether the corresponding $E_{2}$-collapse
is forced by semisimplified Kazhdan--Lusztig Koszulness.

Theorem~\ref*{V1-thm:adams-analog-construction} of Vol~I supplies the
construction. The target functor is built from periodic-CDG
screening-adjoints on the semisimplified KL block
(Section~\ref{sec:part-viii-periodic-cdg}); the analogue of the
mod-$p$ Steenrod algebra is the algebra of periodic-CDG screening
primaries acting on $\mathrm{Ext}^{\bullet}_{\mathrm{KL}}(k,k)$. The
$E_{2}$-collapse follows from Koszul-compatibility of the
semisimplification, pulled back through the bar--cobar adjunction at
admissible level
(Theorem~\ref*{V1-thm:admissible-kl-bar-cobar-adjunction}).

The analogy is now a theorem about screening-primary operations on bar
cohomology. The Adams-analog $E_{2}$-collapse is a legitimate
computational tool on admissible KL blocks.

%% =====================================
\section{Chain-level Chiral Deligne--Tamarkin (associator-dependent)}
\label{sec:part-viii-deligne-tamarkin}

Chain-level chiral Deligne--Tamarkin asks whether the $E_{2}$-action
on $\Zderch(\mathcal{A})$ guaranteed at the level of cohomology lifts
to a strict chain-level $E_{2}$-action. In the classical (non-chiral)
setting, Tamarkin's lift is an $\infty$-quasi-isomorphism depending on
a choice of Drinfeld associator; the chiral theorem must specify its
associator gauge.

Theorem~\ref{thm:chd-deligne-tamarkin} of
Chapter~\ref{ch:chiral-higher-deligne} realises the chiral analogue
in its associator-dependent form. Fixing a Drinfeld associator
$\Phi \in \mathrm{GRT}_{1}(\Q)$, the chain-level $E_{2}$-action on
$\Zderch(\mathcal{A})$ exists as a strict dg-operadic action of
$\mathrm{E}_{2}^{\mathrm{chain}}[\Phi]$, transported from classical
Deligne--Tamarkin through the chiral/algebraic Hochschild model
equivalence of Proposition~\ref{prop:chd-models-equivalent}.
The associator gauge is flagged at
Remark~\ref{rem:chd-associator-discipline}.

The associator-independent formulation is a formal question
(cf.\ Section~\ref{sec:chd-not-proved}): independence is equivalent
to triviality of the $\mathrm{GRT}_{1}(\Q)$-action on the
$E_{2}$-chain-action, and the chiral structure exploits the
non-triviality of the torsor to distinguish faces
(Section~\ref{sec:part-viii-koszulness-moduli}). The
associator-dependent formulation is the appropriate notion for the
fourteen-chart atlas on $M_{\mathrm{Kosz}}$.

%% =====================================
\section{Higher Massey odd-weight obstruction criterion}
\label{sec:part-viii-massey}

Strict Yangian associator coordinates are first detected at length
three. For a finite-dimensional complex simple Lie algebra
$\mathfrak g$, rational Yangian $r$-matrix \(r(u)=\Omega/u\), and
chiral endomorphism operad
\(\mathcal E=\mathrm{End}^{\mathrm{ch}}_{Y_\hbar(\mathfrak g)}\),
Proposition~\ref{prop:massey-rrr-yangian-scope} defines the triple
Massey coset in
\[
H^3(\mathcal E)\Big/
\bigl(r\cdot H^\bullet(\mathcal E)+H^\bullet(\mathcal E)\cdot r\bigr).
\]
The Drinfeld--Kohno depth-three word
\(\theta_3=[\Omega_{12},[\Omega_{13},\Omega_{23}]]\) is the test
vector in this quotient.

For a chosen Drinfeld associator \(\Phi\), the first strict
Yangian associator coordinate is
\[
\omega_3(\Phi)=\sigma_3(\Phi)[\theta_3],
\]
where \(\sigma_3(\Phi)\) is extracted from the length-three
Drinfeld--Kohno coefficient after substituting
\(t_{ij}\mapsto\Omega_{ij}\). The quotient criterion is exact:
\(\omega_3(\Phi)=0\) precisely when either
\(\sigma_3(\Phi)=0\) in the coefficient field or
\(\theta_3\) dies in the specialized chiral endomorphism quotient.
Proposition~\ref{prop:yangian-associator-order-3} identifies the KZ
length-three coefficients, in the standard \(1/(2\pi i)\)
normalization, as lying in
\(\Q\cdot\zeta(3)/(2\pi i)^3\).

Brown's motivic-depth theorem (\cite{BrownMixedTate}) supplies the
odd-weight motivic coordinates of the $\mathrm{GRT}_{1}$ torsor.
An odd-weight coordinate produces a genuine transition between strict
charts of \(M_{\mathrm{Kosz}}(\mathcal A)\) precisely on those chiral
endomorphism charts where its \(\sigma_{2m+1}\)-coefficient is nonzero
and the corresponding Drinfeld--Kohno Lie word survives the Massey
indeterminacy quotient.

%% =====================================
\section{Koszulness Moduli scheme $M_{\mathrm{Kosz}}$ as $\mathrm{GRT}_{1}$-torsor atlas}
\label{sec:part-viii-koszulness-moduli}

The Platonic form of the four theorem threads is the moduli scheme
$M_{\mathrm{Kosz}}(\mathcal{A})$ of Chapter~\ref{chap:koszulness-moduli},
on which the fourteen classical Koszulness characterisations of an
$E_{1}$-chiral algebra $\mathcal{A}$ are affine charts and
$\mathrm{GRT}_{1}(\Q)$ acts by transporting Drinfeld-associator
coordinates. The moduli scheme is Definition~\ref{def:mkosz}; its
atlas structure Theorem~\ref{thm:mkosz-fourteen-charts}; its
shadow-class stratification Theorem~\ref{thm:mkosz-shadow-stratification}.

The fourteen characterisations are local coordinate statements. On
its own chart each is a theorem; on a different chart the additional
conditions are the associator-change cocycles governing transition
functions. The $\mathrm{GRT}_{1}(\Q)$-torsor structure is exhibited in
Section~\ref{sec:mkosz-grt-torsor}.

Sections~\ref{sec:part-viii-curved-dunn},
\ref{sec:part-viii-ds-hochschild},
\ref{sec:part-viii-periodic-cdg}, and
\ref{sec:part-viii-deligne-tamarkin} are coordinate-chart
identifications inside this atlas. In the Kontsevich chart
$\Phi = \Phi_{\mathrm{Kon}}$ (chart $U_{14}$) the chain-level
$E_{2}$-action is explicit; in the Knizhnik--Zamolodchikov chart
$\Phi = \Phi_{\mathrm{KZ}}$ (charts $U_{1}$--$U_{10}$) the admissible
bar--cobar adjunction is explicit. The transition cocycles between
charts are governed by the Yangian obstruction criterion of
Section~\ref{sec:part-viii-massey}.

%% =====================================
\section{Infinite Fingerprint Classification and the G/L/C/M/FF projection}
\label{sec:part-viii-fingerprint}

Chapter~\ref{chap:infinite-fingerprint} is the companion Platonic-form
chapter: classification of $E_{1}$-chiral algebras on the
$C_{2}$-cofinite standard landscape by the six-slot fingerprint
$\varphi'(\mathcal{A}) = (p_{\max}, r_{\max}, \chi_{\mathrm{VOA}},
n_{\mathrm{strong}}, \mathrm{coset}, \kappaChHodge)$.

Slots are defined in Definition~\ref{def:fingerprint-slots}. The
five-slot truncation fails on the Heisenberg family
(Proposition~\ref{prop:five-slot-insufficient-ch}); the sixth slot
$\kappaChHodge$ restores completeness
(Proposition~\ref{prop:six-slot-distinguishes-ch}). The completeness
theorem is Theorem~\ref{thm:fingerprint-complete-ch}: equality of
fingerprints forces isomorphism of chiral algebras.

The G/L/C/M quadrichotomy is the coarse projection of $\varphi'$ onto
$r_{\max}$ on the non-degenerate locus $\kappaChHodge \ne 0$
(Corollary~\ref{cor:double-coarse-projection}). At $\kappaChHodge = 0$
a fifth class FF (Feigin--Frenkel) appears as a genuine new stratum,
populated by the critical-level Feigin--Frenkel centre of
$V_{-h^{\vee}}(\mathfrak{g})$ and related degenerate-sector examples.

The sixth slot $\kappaChHodge$ is the pullback of Vol~I's Theorem~D
universal Arakelov class under a landscape-classifying morphism
(Remark~\ref{rem:fingerprint-theoremD-coupling}).

%% =====================================
\section{Bismut--Gillet--Soulé Quillen norm as BV one-loop determinant}
\label{sec:part-viii-quillen-bv}

The $\SCchtop$ bar-differential on $\mathbf{H}_{\Delta_5}$-observables
produces a factorisation-algebraic partition function $Z^{(1)}$ on
$\mathrm{Ran}(E^{\mathrm{nod,sm}}_{24})$ whose Quillen norm is the
one-loop BV determinant of the twisted 11D supergravity theory on
$K3 \times T^2$:
\[
\log Z^{(1)}(\mathbf{H}_{\Delta_5})
\;=\;
-\log\|\Delta_5\|_Q^2
\;=\;
-\log\Delta_5
\;-\;
24\,L'(0,\Delta_5,\mathrm{std})
\]
combines three inputs on the non-degenerate Humbert-complement locus:
Bismut--Gillet--Soulé 1988 Quillen-norm immersion along the Kuga--Satake
embedding
$\overline{\mathcal{A}_2}\hookrightarrow\overline{\mathcal{A}_{21}}$;
Yoshikawa 2004 equating the analytic torsion with a period of the
Igusa cusp form $\Delta_5$; Bruinier--K\"uhn 2003 equating the boundary
Borcherds-lift contribution along the Humbert wall $H_4$ with
$24\,L'(0,\Delta_5,\mathrm{std})$.

The coefficient $24$ is doubly determined: as the second Chern class
of the rank-$21$ K3 transcendental lattice variation over
$\overline{\mathcal{A}_2}$; as the count of M5-brane sources at
Kodaira $I_1$ fibres on elliptic K3. The Quillen norm $\|\Delta_5\|^2$
is the sole modular-invariant scalar surviving the averaging map
$\mathrm{av}$ applied to $\mathbf{H}_{\Delta_5}$. Under its comparison
hypotheses, the curved-Dunn $H^2$-vanishing of
Section~\ref{sec:part-viii-curved-dunn} supplies the chain-level
scaffold making this a forced one-loop identification rather than a
coincidence between analytic torsion and modular period.

On the K3 pair, Volume~III's universal trace identity (Vol~I--Vol~III
bridge via $\Phi_{\mathrm{hol}}$ and $\Phi$) reads
\[
K(\mathbf{H}_{\Delta_5})
\;=\;
\kappa_{\mathrm{BKM}}(\mathfrak{g}_{\Delta_5})
\;=\;
c_{\Lambda^{3,2}}(0)/2
\;=\;
5
\;=\;
\mathrm{wt}(\Delta_5);
\]
the Quillen-norm $\log Z^{(1)}$ is the differential-geometric companion
whose residue along $H_4$ recovers the same $\Delta_5$ through a
distinct analytic-torsion route.

%% =====================================
\section{Class-$\mathcal{S}$ $\mathcal{T}[A_1,\Sigma_{0,24}]$ Schur-index
cross-check and the Humbert=AD identification}
\label{sec:part-viii-schur-humbert}

The class-$\mathcal{S}$ theory
$\mathcal{T}[A_1, \Sigma_{0,24}]$ constructed by gluing $22$
$A_1$-trinions along $21$ tubes on the $24$-punctured sphere
$\Sigma_{0,24}$ has vector- and hypermultiplet counts
\[
(n_v, n_h)
\;=\;
(63, 88),
\]
by the Chacaltana--Distler 2010 tinkertoy assembly combined with
Shapere--Tachikawa 2008 vector/hypermultiplet counting. Applying the
Beem--Lemos--Liendo--Peelaers--Rastelli--van Rees 2013 4D/2D map,
\[
c_{4d}
\;=\;
\frac{2n_v + n_h}{12}
\;=\;
\frac{126+88}{12}
\;=\;
\frac{107}{6},
\qquad
c_{2d}
\;=\;
-12\,c_{4d}
\;=\;
-214
\;=\;
-2\cdot 107.
\]
The $\chi_{2d}$ chiral algebra of $\mathcal{T}[A_1,\Sigma_{0,24}]$ at
$c = -214$ is a distinct VOA from $\mathbf{H}_{\Delta_5}$. The
integer ``$24$'' — twenty-four regular punctures of $\Sigma_{0,24}$
on $\mathbb{P}^1$, twenty-four Kodaira $I_1$ fibres on elliptic K3 —
reflects a shared rank-$22$ automorphism pattern:
$\operatorname{Aut}(\Sigma_{0,24})\supset S_{24}$;
$\operatorname{NS}(K3)_{\mathrm{prim}}\supset A_1^{22}$. The two chiral
algebras decouple as $\SCchtop$-data:
$\mathrm{av}(\chi_{2d}[\mathcal{T}[A_1,\Sigma_{0,24}]])$ carries
$\kappaChHodge = c_{2d}/2 = -107$, while
$\mathrm{av}(\mathbf{H}_{\Delta_5})$ carries $\kappa_{\mathrm{BKM}} = 5$.

\paragraph{Humbert equals Argyres--Douglas.}
At a generic point $E_{\tau_1}\times E_{\tau_2}$ on the Humbert
divisor $H_4 \subset \overline{\mathcal{A}_2}$, the Seiberg--Witten
curve of the Coulomb branch adjacent to $\mathcal{T}[A_1,\Sigma_{0,24}]$
degenerates to a pair of elliptic curves glued at a nodal point.
Gaiotto--Moore--Neitzke 2009 Example~8.3 identifies this degeneration
with an $(A_1, A_{2N-1})$-type Argyres--Douglas point: the
regular-singular moduli-space direction collapses onto the
irregular-singular direction of an Argyres--Douglas singularity. The
four-fold sibling stratification
$\{\Psi, \Psi^{\deg}, \Psi^{\mathrm{tor}}, \Psi^{\mathrm{metap}}\}$ of
$\mathbf{H}_{\Delta_5}$ coincides with the class-$\mathcal{S}$ join
\[
\{\Psi,\Psi^{\deg},\Psi^{\mathrm{tor}},\Psi^{\mathrm{metap}}\}
\;=\;
\mathrm{CDT}\cup\mathrm{AD}
\;=\;
\{\text{regular},\text{irregular},\text{Argyres--Douglas},\text{twisted}\}.
\]
Three decompositions agree on the single object
$\mathbf{H}_{\Delta_5}$: the Hall--Drinfeld decomposition (the
$\SCchtop$ bar-differential of
Section~\ref{sec:part-viii-quillen-bv}); the K3-lattice decomposition
(Kuga--Satake $K3 \hookrightarrow \operatorname{Sh}(\mathrm{SO}(2,19))$);
the class-$\mathcal{S}$ decomposition (Seiberg--Witten curve
degeneration).

\paragraph{Three-loop Arnold cancellation on $\mathrm{Conf}_4(E)$.}
The four-point function $\phi^{(4)}$ on $\mathrm{Conf}_4(E_\tau)$
sums $10$ Belokurov--Shavgulidze topologies, including ladder
splitting $T_5$ and theta-graph $T_6$ in its two wheel orientations.
The whole sum cancels pairwise under the Arnold--Cohen flip
$H_1 \leftrightarrow H_4$ via the Arnold relation
$\eta_{ij}\wedge\eta_{jk}+\eta_{jk}\wedge\eta_{ki}+\eta_{ki}\wedge\eta_{ij} = 0$
with $\eta_{ij} = d\log\vartheta_1(z_i - z_j,\tau)$; hence
$\phi^{(4)} = 0$ on $\mathrm{Conf}_4(E_\tau)$ for every $\tau$.
Cross-verification via Humbert non-admissibility: at a generic point
of $H_4$ the period degenerates to a pair of elliptic periods and the
admissible $E_8$-weight condition fails by half-integer excess
$\widetilde{D}_4 = 1/2$, which forces $\phi^{(4)}$ to vanish on the
admissible lattice independently. Both routes close the Arnold
cancellation at three loops without appeal to the regularised
Bruinier--Kühn boundary term.

%% =====================================
\section{$W_\infty[c=-214]$ cocycles and the Conway $V^{s\natural}$
metaplectic image}
\label{sec:part-viii-winf-conway}

The topologisation ladder from critical-level class~$\mathbf{M}$
reaches $W_\infty^{\mathrm{top}}$ at its endpoint
(Theorem~\ref{thm:climax-w-infty-endpoint},
Chapter~\ref{ch:programme-climax-platonic}). At the class-$\mathcal{S}$
value $c = -214$ from $\mathcal{T}[A_1,\Sigma_{0,24}]$
(Section~\ref{sec:part-viii-schur-humbert}), the chiral-Hochschild
generators $e_3, e_4, e_5, e_6 \in \ChirHoch^2(\mathrm{Vir}_{-214})$
align explicitly with Pope--Romans--Shen 1990 $W_\infty$-primaries.
The weight-$5$ identification $e_5 = W_5$ is Wang 1998
Proposition~4.2, read on the chiral-Hochschild side; the weight-$4$
identification reads
\[
e_4
\;=\;
W_4 \;-\; \frac{107}{11}\,\Lambda_Z,
\qquad
\Lambda_Z
\;=\;
{:}TT{:}\;-\;\tfrac{3}{10}\partial^2 T,
\]
with the coefficient $107/11$ pinned by Virasoro-primarity of $e_4$
at $c = -214$. Primarity is an open condition on a one-parameter
family; Pope--Romans--Shen picks out the single admissible value. The
sequence $e_3, e_4, e_5, e_6$ realises the Pope--Romans--Shen
$W_\infty$-primary tower as the chiral-Hochschild chart on the
class-$\mathcal{S}$ fixed point.

\paragraph{Conway $V^{s\natural}$ as metaplectic $\Psi^s$-image on the
Leech lattice.}
The Conway module $V^{s\natural}$ of Duncan 2007 (Duke~139) is the
$\mathbb{Z}/2$-orbifold of the $24$-fermion superalgebra
$A(\Lambda_{24})$ on the Leech lattice $\Lambda_{24}$: an $N = 1$
superconformal VOA of central charge $c = 12$. The super-Borcherds
extension of Scheithauer 2008 (Invent.\ Math.\ 172) realises
$V^{s\natural}$ as the $\Psi^s$ image of $\Lambda_{24}$ under the
fermionic vertex superalgebra functor, distinct from the Monster
$\Psi$-image producing $V^\natural$ from the $24$-boson construction
at $c = 24$. The factor of two between the central charges is the
super-bosonisation ratio.

The topologisation ladder has two Leech-lattice fixed points:
$c = 12$ via $V^{s\natural}$; $c = 24$ via $V^\natural$. Together with
the Schellekens row partition
$71 = 24_{\mathrm{Niemeier}} + 1_{V_1=0} + 46_{\mathrm{nonlat}}$
(Theorem~\ref{thm:schellekens-71-all-alpha-zero},
Chapter~\ref{ch:schellekens-71-alpha-classification-platonic}), the
metaplectic companion $V^{s\natural}$ adds a super-sector Leech fixed
point outside the $C_2$-cofinite bosonic standard landscape census.
The $W_\infty[c = -214]$ identification, the
$\mathcal{T}[A_1,\Sigma_{0,24}]$ Schur-index cross-check, and the
$V^{s\natural}$ metaplectic companion are three coordinate readings on
one class-$\mathbf{M}$ topologisation ladder, each specialising at a
distinct central-charge fixed point.

%% =====================================
\section{Concordance: the Platonic form}
\label{sec:part-viii-concordance}

The Part~VIII theorem anchors:

\begin{center}
\renewcommand{\arraystretch}{1.18}
\begin{tabular}{@{}p{0.29\textwidth}p{0.22\textwidth}p{0.40\textwidth}@{}}
\toprule
\textbf{Structure} & \textbf{Mathematical role} & \textbf{Anchor} \\
\midrule
Curved-Dunn $H^{2}=0$, $g \ge 2$ & Conditional theorem &
 Theorem~\ref{thm:curved-dunn-H2-vanishing-all-genera} \\
Modular-bootstrap$\to$curved-Dunn bridge & Proposition &
 Proposition~\ref{prop:modular-bootstrap-to-curved-dunn-bridge} \\
Genus-1 twisted tensor criterion & Proposition &
 Proposition~\ref{prop:genus1-twisted-tensor-product} \\
Irregular-singular KZB regularity from boundary formal type & Conditional theorem &
 Theorem~\ref{thm:irregular-singular-kzb-regularity} \\
Modular operad composition, all $(g,n)$ & Conditional corollary &
 Corollary~\ref{cor:modular-operad-composition-all-levels} \\
F6 genus-three $\lambda_g$ residue & Computed proposition &
 Proposition~\ref{prop:f6-genus-three-lambda-residue} \\
F6 scalar projection versus socle lift & Obstruction proposition &
 Proposition~\ref{prop:f6-scalar-projection-no-multichannel-lift} \\
DS--Hochschild compatibility, class $\mathbf{M}$ & Theorem &
 Theorem~\ref{thm:chd-ds-hochschild} \\
Universal holography, class $\mathbf{M}$ & Corollary &
 Corollary~\ref{cor:universal-holography-class-M} \\
Weight-completed BV--bar, class $\mathbf{M}$ & Proposition &
 Proposition~\ref{prop:bv-bar-class-m-weight-completed} \\
Periodic-CDG Koszul-compatible (admissible) & Theorem (Vol~I) &
 Theorem~\ref*{V1-thm:periodic-cdg-is-koszul-compatible} \\
KL bar--cobar adjunction, admissible level & Theorem (Vol~I) &
 Theorem~\ref*{V1-thm:admissible-kl-bar-cobar-adjunction} \\
Adams-analog construction & Theorem (Vol~I) &
 Theorem~\ref*{V1-thm:adams-analog-construction} \\
Chain-level chiral Deligne--Tamarkin & Theorem (assoc.-dep.) &
 Theorem~\ref{thm:chd-deligne-tamarkin} \\
Yangian Massey obstruction quotient and KZ length-three coefficient &
 Propositions &
 Propositions~\ref{prop:massey-rrr-yangian-scope},
 \ref{prop:yangian-associator-order-3} \\
$M_{\mathrm{Kosz}}(\mathcal{A})$ atlas & Theorem &
 Theorem~\ref{thm:mkosz-fourteen-charts} \\
$M_{\mathrm{Kosz}}$ shadow-class stratification & Theorem &
 Theorem~\ref{thm:mkosz-shadow-stratification} \\
Six-slot fingerprint completeness & Theorem &
 Theorem~\ref{thm:fingerprint-complete-ch} \\
G/L/C/M $\hookleftarrow$ six-slot projection & Corollary &
 Corollary~\ref{cor:double-coarse-projection} \\
BGS Quillen norm as 1-loop BV on $\mathbf{H}_{\Delta_5}$ &
 Identification &
 \S\ref{sec:part-viii-quillen-bv} \\
$\mathcal{T}[A_1,\Sigma_{0,24}]$ Schur index ($c_{2d} = -214$) &
 Cross-check &
 \S\ref{sec:part-viii-schur-humbert} \\
Humbert $H_4$ = Argyres--Douglas $(A_1, A_{2N-1})$ &
 Identification &
 \S\ref{sec:part-viii-schur-humbert} \\
$\phi^{(4)} = 0$ on $\mathrm{Conf}_4(E_\tau)$ via Arnold &
 Identification &
 \S\ref{sec:part-viii-schur-humbert} \\
$W_\infty[c=-214]$ chiral-Hochschild primaries $e_3,\ldots,e_6$ &
 Identification &
 \S\ref{sec:part-viii-winf-conway} \\
Conway $V^{s\natural}$ at $c = 12$ as $\Psi^s(\Lambda_{24})$ &
 Identification &
 \S\ref{sec:part-viii-winf-conway} \\
\bottomrule
\end{tabular}
\end{center}

On the non-degenerate locus of the standard landscape, each row is
read with the hypotheses of its home theorem. Three questions lie
outside the atlas: the degenerate admissible logarithmic lane; the
associator-independent chain-level chiral Deligne--Tamarkin action;
finite improvements to the weight-completed bridge. These are
formulated in their home locations
(Remark~\ref{rem:chd-native-vs-derived},
Section~\ref{sec:chd-not-proved},
Remark~\ref{rem:bv-bar-class-m-frontier}); they are not hypotheses for
the theorem atlas.

The Platonic form is the closed atlas: seven theorems, an infinite
fingerprint, a $\mathrm{GRT}_{1}(\Q)$-parametrised face family,
universal holography extended to class~$\mathbf{M}$, and the
mechanically navigable atlas linking the four theorem threads. The
four frontiers are not independent.

%% =====================================================================
